% Generated draft: requires mathematical and source review.
\documentclass[11pt]{article}
\usepackage[margin=25mm]{geometry}
\usepackage{fontspec}
\setmainfont{Latin Modern Roman}
\usepackage{amsmath,amssymb,mathtools,unicode-math}
\setmathfont{Latin Modern Math}
\usepackage{enumitem,needspace,xurl,hyperref}
\setlength{\parindent}{0pt}
\setlength{\parskip}{5pt}
\setlength{\emergencystretch}{3em}
\newcommand{\N}{\mathbb{N}}
\newcommand{\Z}{\mathbb{Z}}
\newcommand{\Q}{\mathbb{Q}}
\newcommand{\R}{\mathbb{R}}
\newcommand{\C}{\mathbb{C}}
\newcommand{\F}{\mathbb{F}}
\newcommand{\A}{\mathbb{A}}
\newcommand{\PP}{\mathbb{P}}
\newcommand{\E}{\mathbb{E}}
\newcommand{\eps}{\varepsilon}
\newcommand{\dd}{\,\mathrm{d}}
\begin{document}
\section*{UnsolvedMath ID: 800177}
\textbf{Source problem code:} catalogue entry
\subsection{Definitions and mathematical statement}
Let \(\zeta(s)\) denote the Riemann zeta function, defined for \(\operatorname{Re}s>1\) by \(\sum_{n=1}^{\infty}n^{-s}\) and elsewhere by meromorphic continuation. The Riemann hypothesis (RH) asserts that every zero \(\rho\) with \(0<\operatorname{Re}\rho<1\) satisfies \(\operatorname{Re}\rho=\tfrac12\). All logarithms below are natural.

The source writes \(S(t)=\pi^{-1}\arg\zeta(\tfrac12+it)\) without specifying the argument convention or values at zero ordinates. For a precise candidate formulation under RH, for \(t>0\) that is not a zero ordinate, define the argument by continuous variation along the line segments from \(2\) to \(2+it\) and then to \(\tfrac12+it\), starting with \(\arg\zeta(2)=0\). Set
\[
S(t)=\frac{1}{\pi}\arg\zeta\!\left(\tfrac12+it\right).
\]
At a positive zero ordinate \(\gamma\), use the convention
\[
S(\gamma)=\frac12\left(\lim_{h\downarrow0}S(\gamma-h)+\lim_{h\downarrow0}S(\gamma+h)\right).
\]
These argument and boundary conventions are explicit interpretations for this draft and require confirmation against the source.

The supplied record reports that Carneiro, Chandee, and Milinovich [CCM] proved, assuming RH,
\[
|S(t)|\leq\left(\frac14+o(1)\right)\frac{\log t}{\log\log t}
\qquad(t\to+\infty).
\]
\textbf{Requested improvement.} Assuming RH, establish the existence of a fixed real constant \(C\) with \(0\leq C<\tfrac14\) such that, for every \(\varepsilon>0\), there is a real number \(T_{\varepsilon}>e\) for which
\[
|S(t)|\leq(C+\varepsilon)\frac{\log t}{\log\log t}
\qquad\text{for every real }t\geq T_{\varepsilon}.
\]
Thus the requested reduction concerns the leading constant in the asymptotic bound. The source does not specify a particular smaller target constant.
\subsection{Short English statement}
Assuming the Riemann hypothesis, improve the supplied asymptotic bound for the normalized argument of the Riemann zeta function by replacing its leading constant \(1/4\) with a fixed smaller constant. The argument normalization and treatment of zero ordinates have been made explicit as draft conventions requiring review.
\subsection{Sources}
\begin{itemize}
\item Local UnsolvedMath registry, permanent ID \(800177\), exact source problem code \texttt{catalogue entry}; title: \emph{Moments of L-functions --- Moments of the Riemann zeta function --- Problem 28}. The registry credits Hamieh, Kadiri, and Ng, with proposed year 2024, and labels the problem \texttt{open}; this is a supplied status, not an independently verified assessment.
\item Hamieh, Kadiri, and Ng, \emph{Moments of L-functions Problem List} (2024), Problem 28, \url{https://arxiv.org/abs/2405.17800}. The supplied item attributes the question to Nathan Ng, Lethbridge. Supplied provenance records access on 2026-10-16 and extraction by \texttt{arxiv\_latex}; these are source-record annotations, not a claim of independent access or verification for this draft.
\item Carneiro, Chandee, and Milinovich, reference label [CCM] in the supplied statement, are credited there with the conditional bound having leading constant \(1/4\). No full bibliographic entry or separate URL for [CCM] is supplied.
\end{itemize}
\end{document}
